% !TeX root = ../main.tex
\chapter{Resultados}\label{cap:resultados}

Empieza describiendo el protocolo experimental (datos, métricas,
semillas…) para que los resultados sean reproducibles.

\section{Protocolo experimental}

\section{Resultados principales}

Las figuras con datos se generan con los scripts de
\texttt{scripts/figuras/} (\cref{fig:ejemplo}): nunca a mano ni como
capturas de pantalla. Así usan la misma tipografía que la memoria, llevan
coma decimal y se pueden regenerar si cambian los datos.

\begin{figure}[H]
  \centering
  \includegraphics[width=0.8\linewidth]{generadas/ejemplo_curva}
  \caption{Figura de ejemplo generada con \texttt{scripts/figuras/ejemplo\_curva.py}.}
  \label{fig:ejemplo}
\end{figure}

Las tablas de resultados también pueden generarse desde el código
(\cref{tab:ejemplo}), lo que evita errores al copiar cifras a mano.

% Generado automáticamente por scripts/figuras. NO EDITAR A MANO.
\begin{table}[H]
  \centering
  \caption{Métricas de error del ejemplo (generada desde Python).}
  \label{tab:ejemplo}
  \begin{tabular}{lc}
    \toprule
    Métrica & Valor \\
    \midrule
    MAE & \num{0.021} \\
    RMSE & \num{0.026} \\
    Error máximo & \num{0.087} \\
    $\tau$ (s) & \num{2.500} \\
    \bottomrule
  \end{tabular}
\end{table}


\section{Discusión}

Interpreta los resultados: qué significan, por qué salen así y qué
limitaciones tienen.
